%=============================================================================
\section{Demographic Positional Encoding}
\label{sec:dpe}
%=============================================================================

Having established that each model carries a distinct, internally consistent
bias fingerprint (\Cref{sec:results}), we ask whether these disparities can be
\emph{reduced at inference time} without retraining, fine-tuning, or access to
model weights beyond a forward pass. We introduce \textbf{Demographic Positional
Encoding} (DPE), a training-free intervention that estimates a model's
per-group bias offset from its own baseline behaviour and counteracts it by
injecting a demographic-conditioned corrective signal into the model's visual
token embeddings.

The method rests on a simple hypothesis: if a model systematically shifts its
outputs for a demographic group $g$ relative to the population mean, that shift
corresponds to a displacement in the model's internal representation space, and
an equal and opposite displacement---applied to the visual tokens before the
language decoder attends to them---will move the group's outputs back toward
parity. DPE operationalises this in three stages: (i) bias-vector estimation,
(ii) projection into embedding space, and (iii) inference-time injection.

%-----------------------------------------------------------------------------
\subsection{Notation and problem setup}
%-----------------------------------------------------------------------------

Let $\mathcal{D} = \{(x_i, g_i)\}_{i=1}^{N}$ be the evaluation set, where $x_i$
is an image and $g_i \in \mathcal{G}$ its demographic group. We define
$\mathcal{G}$ as the Cartesian product of gender presentation and jurisdictional
region, $\mathcal{G} = \mathrm{Gender} \times \mathrm{Region}$, yielding the
groups over which disparity is measured in \Cref{sec:results}. For each
image--probe response we obtain a vector of $K=3$ scored bias dimensions,
\begin{equation}
    s_i = \big(\,\mathrm{valence}_i,\;\mathrm{stereotype}_i,\;\mathrm{confidence}_i\,\big) \in \mathbb{R}^{3},
\end{equation}
produced by the LLM judge (\Cref{sec:scoring}). A VLM exposes a vision encoder
$f_\text{vis}$ that maps an image to a sequence of $T$ visual token embeddings
$H = f_\text{vis}(x) \in \mathbb{R}^{T \times d}$, where $d$ is the hidden width
of the representation consumed by the language decoder.

%-----------------------------------------------------------------------------
\subsection{Stage 1: Demographic bias-vector estimation}
%-----------------------------------------------------------------------------

From the model's baseline responses we compute the grand mean score vector and
the per-group mean,
\begin{equation}
    \mu = \frac{1}{N}\sum_{i=1}^{N} s_i,
    \qquad
    \mu_g = \frac{1}{|\mathcal{I}_g|}\sum_{i \in \mathcal{I}_g} s_i,
\end{equation}
where $\mathcal{I}_g = \{i : g_i = g\}$ indexes the responses for group $g$. The
\emph{bias offset} of group $g$ is the signed deviation of its mean from the
population mean, and the corresponding \emph{correction vector} is its negation,
\begin{equation}
    \delta_g \;=\; \mu - \mu_g \;\in\; \mathbb{R}^{3}.
    \label{eq:correction}
\end{equation}
A group treated more negatively than the population on some dimension (e.g.\ a
lower mean valence) receives a positive correction on that dimension, and vice
versa; a perfectly equitable group has $\delta_g = \mathbf{0}$. Groups with
fewer than $n_\text{min}=30$ scored responses are excluded from estimation to
avoid unstable offsets, and fall back to $\delta_g = \mathbf{0}$ (no
intervention). Crucially, $\delta_g$ is estimated from the target model's
\emph{own} baseline outputs, making DPE self-calibrating and model-specific: the
same procedure yields a different correction field for each VLM.

%-----------------------------------------------------------------------------
\subsection{Stage 2: Projection into embedding space}
%-----------------------------------------------------------------------------

The correction $\delta_g$ lives in the $3$-dimensional score space, whereas the
intervention must act in the $d$-dimensional visual embedding space
($d \gg 3$). We map between them with a fixed linear operator
$W \in \mathbb{R}^{d \times 3}$ and define the group's positional encoding
\begin{equation}
    \varepsilon_g \;=\; W\,\delta_g \;\in\; \mathbb{R}^{d}.
    \label{eq:encoding}
\end{equation}
We construct $W$ to have orthonormal columns by QR-decomposing a Gaussian random
matrix, $W = \mathrm{qr}(G)$ with $G_{jk} \sim \mathcal{N}(0,1)$, under a fixed
random seed. Orthonormality ensures the projection is an isometry onto its
range, $\lVert \varepsilon_g \rVert_2 = \lVert \delta_g \rVert_2$, so the
\emph{magnitude} of a group's bias offset is preserved as the magnitude of its
representational shift, and distinct score-space directions map to distinct,
non-interfering embedding-space directions. Fixing the seed makes $W$
deterministic and reproducible across runs and models, so the encoding for a
given $\delta_g$ is identical every time. We deliberately avoid learning $W$:
a learned projection would require gradient access and a training objective,
reintroducing exactly the retraining cost DPE is designed to avoid. Learning $W$
under a fairness objective is a natural extension (\Cref{sec:limitations}).

%-----------------------------------------------------------------------------
\subsection{Stage 3: Inference-time injection}
%-----------------------------------------------------------------------------

At inference, given an image $x$ of a subject in group $g$, we apply the
encoding additively to every visual token produced by the vision encoder,
\begin{equation}
    \tilde{H} \;=\; H + \alpha\,\mathbf{1}_T\,\varepsilon_g^{\top},
    \qquad
    \tilde{H}_t = H_t + \alpha\,\varepsilon_g \;\; \forall t,
    \label{eq:inject}
\end{equation}
where $\mathbf{1}_T \in \mathbb{R}^{T}$ broadcasts the encoding across all $T$
tokens and $\alpha \ge 0$ is a scalar \emph{correction strength}. The shifted
tokens $\tilde{H}$ replace $H$ in the language decoder's cross-attention, so the
model conditions its generation on a representation that has been nudged toward
the population mean for group $g$. Implementation-wise this requires no weight
modification: we register a forward hook on the vision encoder module that
intercepts $H$ and returns $\tilde{H}$, leaving the rest of the forward pass
untouched. Setting $\alpha = 0$ recovers the unmodified baseline; larger $\alpha$
applies a stronger corrective push. Because $\varepsilon_g$ depends only on the
group label---available from dataset metadata---DPE adds no per-image inference
cost beyond a single vector addition.

\begin{algorithm}[t]
\caption{Demographic Positional Encoding (DPE)}
\label{alg:dpe}
\begin{algorithmic}[1]
\Require baseline scores $\{(s_i, g_i)\}$; VLM with vision encoder $f_\text{vis}$;
         strength $\alpha$; seed $\rho$
\Statex \textbf{// Stage 1: estimate correction field (once)}
\State $\mu \gets \mathrm{mean}_i\, s_i$
\For{each group $g \in \mathcal{G}$ with $|\mathcal{I}_g| \ge n_\text{min}$}
    \State $\delta_g \gets \mu - \mathrm{mean}_{i \in \mathcal{I}_g}\, s_i$
\EndFor
\Statex \textbf{// Stage 2: fixed orthonormal projection}
\State $W \gets \mathrm{qr}(G),\; G \sim \mathcal{N}(0,1)^{d\times 3}$ seeded by $\rho$
\Statex \textbf{// Stage 3: inference with injection}
\For{each image $x$ of group $g$}
    \State $H \gets f_\text{vis}(x)$
    \State $\varepsilon_g \gets W\,\delta_g$
    \State $\tilde{H} \gets H + \alpha\,\varepsilon_g$ \Comment{broadcast over tokens}
    \State generate response from $\tilde{H}$
\EndFor
\end{algorithmic}
\end{algorithm}

%-----------------------------------------------------------------------------
\subsection{Evaluation protocol}
%-----------------------------------------------------------------------------

We evaluate DPE by re-running the full probe battery with injection enabled and
comparing against the untouched baseline on the identical set of images and
probes, so that any change is attributable solely to the intervention. We report
three complementary quantities. First, the \emph{demographic disparity}
$\sigma_A$ for a demographic axis $A$ (gender or region), defined as the
standard deviation of per-group mean valence across the groups of $A$; lower
$\sigma_A$ indicates more equitable treatment, and we summarise the effect as
the relative reduction $(\sigma_A^\text{base} - \sigma_A^\text{DPE}) /
\sigma_A^\text{base}$. Second, per-dimension effect sizes (Cohen's $d$) and
Welch's $t$-test $p$-values between baseline and DPE score distributions, to
establish that observed shifts are statistically reliable and to quantify their
magnitude. Third, a per-region $\times$ per-probe map of the mean valence change
$\Delta = \bar{s}^\text{DPE} - \bar{s}^\text{base}$, to verify that the
intervention moves the most disadvantaged groups upward rather than merely
compressing variance by penalising advantaged groups. A successful intervention
reduces $\sigma_A$ while raising the valence of the lowest-scoring groups toward
the mean.
